\documentclass[conference]{IEEEtran}

\usepackage[utf8]{inputenc}
\usepackage[T1]{fontenc}
\usepackage[turkish]{babel}
\usepackage{cite}
\usepackage{amsmath,amssymb}
\usepackage{graphicx}
\usepackage{booktabs}
\usepackage{array}
\usepackage{url}
\usepackage{float}
\usepackage{xcolor}
\usepackage{listings}

\graphicspath{{./}}

\definecolor{codegray}{RGB}{245,245,245}
\definecolor{codegreen}{RGB}{0,128,0}
\definecolor{codered}{RGB}{180,0,0}
\definecolor{accentblue}{RGB}{0,119,182}

\lstset{
  backgroundcolor=\color{codegray},
  basicstyle=\ttfamily\footnotesize,
  breaklines=true,
  keywordstyle=\color{accentblue}\bfseries,
  commentstyle=\color{codegreen},
  stringstyle=\color{codered},
  frame=single,
  language=Python
}

\begin{document}

\title{Top1--Top2 Token Olasılık Farkından Özellik Çıkararak\\
Cosmos-T1 Cevap Doğruluğu Tahmini}

\author{
  \IEEEauthorblockN{Betül Arslan}
  \IEEEauthorblockA{
    Hesaplamalı Anlambilim -- Proje Konusu 4\\
    \texttt{betullarslan@gmail.com}
  }
}

\maketitle

\begin{abstract}
Bu çalışmada, büyük dil modellerinin cevap üretimi sırasında oluşan token
düzeyindeki olasılık farklarının, üretilen cevabın doğru olup olmadığını tahmin
etmek için kullanılıp kullanılamayacağı incelenmiştir. Cosmos-T1
(\texttt{ytu-ce-cosmos/Turkish-Gemma-9b-T1}) modeli, Türkçe GSM8K
(\texttt{gsm8k\_tr}) veri kümesindeki matematik soruları üzerinde
çalıştırılmıştır. Model her token ürettiğinde en yüksek olasılıklı token ile
ikinci en yüksek olasılıklı token arasındaki fark kaydedilmiş, daha sonra bu
fark dizilerinden 14 istatistiksel özellik çıkarılmıştır. Toplanan 357 kaydın
187'si doğru, 170'i yanlış etiketlenmiştir. Sınıf dengesi sağlanarak 120 doğru
+ 120 yanlış örnek eğitim, 50 doğru + 50 yanlış örnek test için kullanılmıştır.
Logistic Regression, Random Forest ve Gradient Boosting modelleri
karşılaştırılmış; en dengeli sonuç Gradient Boosting ile elde edilmiştir
(Accuracy=0.72, F1=0.72, ROC-AUC=0.762). Bulgular, top1--top2 fark dizisinin
tek başına kusursuz bir ayırıcı olmadığını; ancak bu diziden çıkarılan
istatistiksel özelliklerin cevap doğruluğu hakkında anlamlı sinyal taşıdığını
göstermektedir.
\end{abstract}

\begin{IEEEkeywords}
büyük dil modelleri, top1-top2 farkı, cevap doğruluğu, Cosmos-T1, GSM8K-TR,
makine öğrenmesi
\end{IEEEkeywords}

\section{Giriş}

Büyük dil modelleri her üretim adımında kelime dağarcığı üzerinde bir olasılık
dağılımı hesaplar ve bu dağılımdan bir sonraki tokeni üretir. Bu dağılım yalnızca
sonraki token seçimi için değil, modelin o adımda ne kadar kararlı olduğunu
anlamak için de kullanılabilir. Bu çalışmada kullanılan temel sinyal,
her üretim adımında en yüksek olasılıklı token ile ikinci en yüksek olasılıklı
token arasındaki farktır.

Bir üretim adımı $t$ için top1--top2 farkı şu şekilde tanımlanmıştır:

\begin{equation}
  \delta_t =
  P(\hat{y}^{(1)}_t \mid x, y_{<t}) -
  P(\hat{y}^{(2)}_t \mid x, y_{<t}),
  \quad t=1,\ldots,T .
  \label{eq:delta}
\end{equation}

Burada $\hat{y}^{(1)}_t$ ve $\hat{y}^{(2)}_t$, sırasıyla $t$. adımda en yüksek
ve ikinci en yüksek olasılığa sahip tokenleri; $x$ giriş sorusunu, $y_{<t}$
önceki tokenleri ve $T$ üretilen toplam token sayısını göstermektedir.
$\delta_t$ değeri 0 ile 1 arasındadır. Değerin yüksek olması modelin ilk token
seçimini ikinci seçeneğe göre daha güçlü tercih ettiğini, düşük olması ise iki
seçenek arasında daha kararsız kaldığını gösterir.

Bu çalışmanın araştırma sorusu şudur: Cosmos-T1 modeli bir matematik sorusuna
cevap üretirken oluşan top1--top2 fark dizisi, cevabın doğru mu yanlış mı
olduğunu tahmin etmek için kullanılabilir mi?

\section{İlgili Çalışmalar}

Sinir ağı çıktılarındaki olasılık değerlerinin güven ölçütü olarak kullanılması
uzun süredir çalışılan bir konudur. Hendrycks ve Gimpel, maksimum softmax
olasılığının yanlış sınıflandırılmış ve dağılım dışı örneklerin tespitinde basit
ama etkili bir taban çizgisi sunduğunu göstermiştir~\cite{hendrycks2016}.
Bu çalışma, modelin çıktı olasılıklarının hata tahmini için bilgi taşıyabileceği
fikrini desteklemektedir.

Büyük dil modelleri özelinde Kadavath ve arkadaşları, modellerin kendi
cevaplarının doğruluğu hakkında belirli ölçüde bilgi taşıyabildiğini
tartışmıştır~\cite{kadavath2022}. SelfCheckGPT çalışması ise dış bilgi kaynağı
kullanmadan model çıktılarını karşılaştırarak güvenilirlik değerlendirmesi
yapılabileceğini göstermiştir~\cite{selfcheckgpt}. Bizim çalışmamız bu çizgide,
dış doğrulayıcı kullanmak yerine üretim sırasında oluşan iç olasılık sinyalini
incelemektedir.

Matematik problemi çözme tarafında GSM8K veri kümesi ve doğrulayıcı modeller
üzerine yapılan çalışma, dil modellerinin matematik muhakemesi çıktılarının
ayrı bir doğrulama aşamasıyla değerlendirilebileceğini göstermiştir~\cite{gsm8k}.
Bu projede benzer biçimde cevap doğruluğu tahmin edilmektedir; ancak doğrulayıcı
model doğrudan metin üzerinden değil, token olasılık farklarından çıkarılan
sayısal özellikler üzerinden eğitilmiştir.

\section{Model ve Veri Kümesi}

\subsection{Cosmos-T1 Modeli}

Deneylerde \texttt{ytu-ce-cosmos/Turkish-Gemma-9b-T1} modeli kullanılmıştır.
Model, Türkçe akıl yürütme ve matematik problemleri için yayımlanmış Cosmos-T1
ailesine ait bir Gemma tabanlı dil modelidir~\cite{cosmos}. Model kartındaki
önerilere uygun olarak \texttt{temperature=0.6}, \texttt{top\_p=0.95},
\texttt{top\_k=20} ve \texttt{do\_sample=True} ayarları kullanılmıştır.
Greedy decoding tercih edilmemiştir.

\subsection{GSM8K-TR Veri Kümesi}

Veri kümesi olarak \texttt{ytu-ce-cosmos/gsm8k\_tr} kullanılmıştır~\cite{gsm8k_tr}.
Bu veri kümesi, GSM8K matematik problemi veri kümesinin Türkçeye uyarlanmış
sürümüdür. Veri kümesi sayfasında belirtildiği gibi sorular DeepL ile Türkçeye
çevrilmiş, cevaplar ise Türkçe sorular üzerinden GPT-4o ile yeniden
üretilmiştir.

Orijinal GSM8K'de cevaplar çoğunlukla \texttt{\#\#\#\# <sayı>} biçiminde
verilirken, bu Türkçe sürümde cevaplar düz metin halindedir. Bu nedenle altın
cevap son satırlardaki sayısal ifadelerden çıkarılmıştır. Saat biçimli cevaplar
(\texttt{10:30} gibi) sayı çıkarımında hataya neden olabileceği için filtrelenmiştir.
Toplam 34 saat biçimli örnek elenmiş ve 8.734 temiz örnek elde edilmiştir.

\section{Yöntem}

\subsection{Top1--Top2 Farklarının Kaydedilmesi}

Üretim sırasında her token için top1--top2 farkını almak için iki yöntem
denenmiştir. İlk yöntemde \texttt{output\_scores=True} kullanılmıştır; ancak bu
yöntem her adımda tüm kelime dağarcığına ait skorları sakladığı için bellek
tüketimini artırmaktadır. Son sürümde yalnızca gerekli skalar farkı saklayan
\texttt{LogitsProcessor} yaklaşımı kullanılmıştır.

\begin{lstlisting}
class Top1Top2Kaydedici(LogitsProcessor):
    def __init__(self):
        self.farklar = []

    def __call__(self, input_ids, scores):
        probs = torch.softmax(scores.float(), dim=-1)
        top2 = torch.topk(probs, k=2, dim=-1).values
        fark = (top2[:, 0] - top2[:, 1]).detach().cpu().tolist()
        self.farklar.append(fark)
        return scores
\end{lstlisting}

Bu yöntemle tüm logit vektörü yerine her adımda yalnızca $\delta_t$ değeri
kaydedilmiştir. Böylece uzun cevaplarda bellek kullanımı daha kontrollü
tutulmuştur.

\subsection{Üretim Ayarları}

\begin{table}[H]
  \centering
  \caption{Model üretim ayarları}
  \label{tab:params}
  \begin{tabular}{ll}
    \toprule
    \textbf{Parametre} & \textbf{Değer} \\
    \midrule
    Model & Turkish-Gemma-9b-T1 \\
    \texttt{max\_new\_tokens} & 1024 \\
    \texttt{temperature} & 0.6 \\
    \texttt{top\_p} & 0.95 \\
    \texttt{top\_k} & 20 \\
    \texttt{do\_sample} & True \\
    Donanım & A100, bfloat16 \\
    \bottomrule
  \end{tabular}
\end{table}

\subsection{Özellik Mühendisliği}

Her cevap için top1--top2 farkları
$\boldsymbol{\delta}=[\delta_1,\delta_2,\ldots,\delta_T]$ dizisi olarak
saklanmıştır. Klasik makine öğrenmesi algoritmaları sabit boyutlu giriş
gerektirdiği için bu değişken uzunluklu dizi 14 özellikten oluşan sabit boyutlu
bir vektöre dönüştürülmüştür.

Kullanılan özellikler şunlardır:

\begin{itemize}
  \item \texttt{n\_tokens}: Üretilen token sayısıdır. Çok uzun üretimler modelin
        çözümü tamamlamakta zorlandığını veya token limitine yaklaştığını
        gösterebilir.
  \item \texttt{mean}: Ortalama top1--top2 farkıdır. Genel güven düzeyini özetler.
  \item \texttt{std}: Fark dizisinin standart sapmasıdır. Güvenin cevap boyunca
        ne kadar dalgalandığını ölçer.
  \item \texttt{min}, \texttt{max}: En düşük ve en yüksek fark değerleridir.
        En kararsız ve en kararlı üretim anlarını temsil eder.
  \item \texttt{median}: Ortanca fark değeridir. Aykırı değerlerden ortalamaya
        göre daha az etkilenir.
  \item \texttt{q25}, \texttt{q75}, \texttt{iqr}: Dağılımın çeyrekleri ve
        çeyrekler arası aralığıdır. Güven değerlerinin yayılımını özetler.
  \item \texttt{mean\_first10}: İlk 10 tokenin ortalama farkıdır. Cevabın
        başlangıcındaki güveni ölçer.
  \item \texttt{mean\_last10}: Son 10 tokenin ortalama farkıdır. Nihai cevaba
        yaklaşırken modelin güvenini temsil eder.
  \item \texttt{prop\_low\_01}, \texttt{prop\_low\_02}: Farkı sırasıyla 0.10 ve
        0.20 altında olan token oranlarıdır. Düşük güvenli token yoğunluğunu
        ölçer.
  \item \texttt{trend}: Fark dizisine uydurulan doğrusal eğimin katsayısıdır.
        Cevap ilerledikçe güvenin artıp azaldığını gösterir.
\end{itemize}

Temel özellik formülleri aşağıda verilmiştir:

\begin{align}
  \mu &= \frac{1}{T}\sum_{t=1}^{T}\delta_t, \\
  \sigma &= \sqrt{\frac{1}{T}\sum_{t=1}^{T}(\delta_t-\mu)^2}, \\
  \mathrm{IQR} &= Q_{75}-Q_{25}, \\
  p_{<\theta} &= \frac{1}{T}\sum_{t=1}^{T}\mathbf{1}[\delta_t < \theta],
  \quad \theta \in \{0.10, 0.20\}.
\end{align}

\subsection{Eğitim ve Test Bölümlemesi}

Toplanan 357 kaydın 187'si doğru, 170'i yanlış cevap olarak etiketlenmiştir.
Sınıf dengesini sağlamak için her iki sınıftan 170'er örnek kullanılmıştır.
Bunların 120'şer tanesi eğitim, 50'şer tanesi test kümesine ayrılmıştır.

\begin{table}[H]
  \centering
  \caption{Dengeli eğitim ve test bölümlemesi}
  \label{tab:split}
  \begin{tabular}{lrrr}
    \toprule
    \textbf{Küme} & \textbf{Doğru} & \textbf{Yanlış} & \textbf{Toplam} \\
    \midrule
    Eğitim & 120 & 120 & 240 \\
    Test & 50 & 50 & 100 \\
    \textbf{Toplam} & \textbf{170} & \textbf{170} & \textbf{340} \\
    \bottomrule
  \end{tabular}
\end{table}

Üç sınıflandırıcı karşılaştırılmıştır: Logistic Regression, Random Forest ve
Gradient Boosting. Logistic Regression için \texttt{StandardScaler}
kullanılmıştır. Random Forest 300 ağaçla eğitilmiş, Gradient Boosting için
Scikit-learn varsayılan ayarları kullanılmıştır.

\section{Deneysel Sonuçlar}

\subsection{Model Performansı}

\begin{table}[H]
  \centering
  \caption{Model karşılaştırması (240 eğitim, 100 test)}
  \label{tab:results}
  \begin{tabular}{lccccc}
    \toprule
    \textbf{Model} & \textbf{Acc} & \textbf{Prec} & \textbf{Rec} & \textbf{F1} & \textbf{AUC} \\
    \midrule
    Logistic Regression & 0.720 & 0.844 & 0.540 & 0.659 & 0.738 \\
    Random Forest       & 0.700 & 0.727 & 0.640 & 0.681 & 0.747 \\
    Gradient Boosting   & 0.720 & 0.720 & 0.720 & 0.720 & 0.762 \\
    \bottomrule
  \end{tabular}
\end{table}

Sonuçlara göre en dengeli performans Gradient Boosting modeliyle elde edilmiştir.
Bu model test kümesinde 0.72 doğruluk, 0.72 F1 ve 0.762 ROC-AUC değerlerine
ulaşmıştır. Logistic Regression yanlış cevapları yakalamada daha temkinli
davranmış, ancak doğru cevapların önemli bir kısmını yanlış sınıfına atamıştır.
Random Forest ise iki sınıf arasında daha dengeli, fakat Gradient Boosting'e göre
biraz daha düşük F1 skoru üretmiştir.

\subsection{Token Bazlı Top1--Top2 Fark Grafiği}

\begin{figure}[H]
  \centering
  \includegraphics[width=\linewidth]{fig_token_level_diff.png}
  \caption{Bir doğru ve bir yanlış örnek için token bazlı top1--top2 fark dizisi.}
  \label{fig:token_diff}
\end{figure}

Şekil~\ref{fig:token_diff}, doğru ve yanlış cevaplanan iki örnekte top1--top2
farkının token boyunca nasıl değiştiğini göstermektedir. Her iki çizgide de
farklar sık sık 1'e yaklaşmaktadır; bu durum modelin birçok token için güçlü
bir ilk tercih yaptığını gösterir. Ancak doğru ve yanlış örneklerin tamamen
ayrışmadığı görülmektedir. Bu nedenle tek bir token düzeyi grafiği yerine
fark dizisinden istatistiksel özellikler çıkarılmıştır.

\subsection{Ortalama Fark Dağılımı}

\begin{figure}[H]
  \centering
  \includegraphics[width=\linewidth]{fig_mean_diff_distribution.png}
  \caption{Doğru ve yanlış cevaplarda ortalama top1--top2 fark dağılımı.}
  \label{fig:dist}
\end{figure}

Şekil~\ref{fig:dist}'te doğru ve yanlış cevapların ortalama fark dağılımları
kısmen örtüşmektedir. Bu sonuç, yalnızca ortalama top1--top2 farkının cevap
doğruluğunu ayırmak için yeterli olmadığını göstermektedir. Buna rağmen dağılımın
tamamen rastgele olmaması, ortalama farkın diğer özelliklerle birlikte faydalı
olabileceğini göstermektedir.

\subsection{ROC Eğrileri}

\begin{figure}[H]
  \centering
  \includegraphics[width=0.95\linewidth]{fig_roc.png}
  \caption{Üç modelin ROC eğrileri. Kesikli çizgi rastgele tahmin düzeyini gösterir.}
  \label{fig:roc}
\end{figure}

ROC eğrileri tüm modellerin rastgele tahmin çizgisinin üzerinde kaldığını
göstermektedir. En yüksek AUC değeri 0.762 ile Gradient Boosting modelinde
elde edilmiştir. Bu sonuç, top1--top2 farklarından çıkarılan özelliklerin cevap
doğruluğu tahmini için anlamlı sinyal taşıdığını desteklemektedir.

\subsection{Confusion Matrix}

\begin{figure}[H]
  \centering
  \includegraphics[width=\linewidth]{fig_confusion_matrix.png}
  \caption{Logistic Regression, Random Forest ve Gradient Boosting modellerinin confusion matrix sonuçları.}
  \label{fig:confusion}
\end{figure}

Confusion matrix sonuçları modellerin hata karakterini göstermektedir.
Logistic Regression 50 yanlış cevabın 45'ini doğru biçimde yanlış olarak
sınıflandırmıştır; ancak 50 doğru cevabın 23'ünü yanlış sanmıştır. Gradient
Boosting ise her iki sınıfta da 36/50 doğru sınıflandırma yaparak en dengeli
sonucu vermiştir.

\subsection{Özellik Önemi}

\begin{figure}[H]
  \centering
  \includegraphics[width=\linewidth]{fig_feature_importance.png}
  \caption{Random Forest ve Gradient Boosting modellerinde özellik önemleri.}
  \label{fig:feat}
\end{figure}

Özellik önemi grafiğinde \texttt{n\_tokens}, \texttt{mean\_first10},
\texttt{prop\_low\_01}, \texttt{trend}, \texttt{q75} ve
\texttt{mean\_last10} öne çıkmaktadır. Bu sonuç, yalnızca genel ortalama güvenin
değil, cevabın uzunluğunun, düşük güvenli token oranının ve üretim boyunca güven
eğiliminin de doğru/yanlış ayrımı için etkili olduğunu göstermektedir.

\section{Tartışma}

Deneylerde Gradient Boosting'in daha dengeli sonuç vermesi, top1--top2 fark
özellikleri arasındaki doğrusal olmayan ilişkilerin önemli olduğunu
düşündürmektedir. Logistic Regression doğrusal bir karar sınırı kullandığından
yanlış cevapları yakalamakta başarılı olsa da doğru cevapları daha fazla
kaçırmıştır. Ağaç tabanlı modeller ise özellikle \texttt{n\_tokens},
\texttt{trend} ve düşük güven oranı gibi özellikler arasındaki etkileşimleri
daha iyi yakalamıştır.

Özellik öneminde \texttt{n\_tokens}'ın güçlü çıkması dikkatle yorumlanmalıdır.
Bu özellik doğrudan güven ölçütü değildir; modelin cevabı ne kadar uzattığını
ve 1024 token sınırına yaklaşıp yaklaşmadığını yansıtır. Cosmos-T1 uzun
\texttt{<think>} blokları ürettiği için bazı yanlış cevaplarda modelin çözümü
tamamlayamadan token limitine yaklaştığı görülmektedir. Dolayısıyla
\texttt{n\_tokens} bir vekil sinyal olarak çalışmaktadır.

\section{Sınırlılıklar}

Bu çalışmanın temel sınırlılığı veri toplama süresidir. Proje gereksiniminde
1000 eğitim ve 100 test örneği hedeflenmiş olsa da, Cosmos-T1 modelinin uzun
akıl yürütme çıktıları ve Colab GPU süresi nedeniyle bu çalışmada 240 eğitim ve
100 test örneğiyle deney yapılmıştır. Test kümesi sınıf dengesi açısından
gereksinimi karşılamaktadır; ancak eğitim kümesi hedeflenen büyüklüğün
altındadır.

İkinci sınırlılık, \texttt{gsm8k\_tr} cevaplarının düz metin biçiminde olmasıdır.
Altın cevaplar düzenli ifadelerle çıkarılmıştır; bu süreçte az sayıda etiket
gürültüsü kalmış olabilir. Üçüncü sınırlılık ise \texttt{max\_new\_tokens=1024}
değeridir. Daha yüksek token limiti, modelin bazı cevapları tamamlamasına izin
verebilir ve \texttt{n\_tokens} özelliğinin baskınlığını azaltabilir.

\section{Sonuç}

Bu çalışmada Cosmos-T1 modelinin Türkçe GSM8K sorularına verdiği cevapların
doğruluğu, cevap üretimi sırasında oluşan top1--top2 token olasılık farklarından
tahmin edilmeye çalışılmıştır. Her cevap için fark dizisinden 14 istatistiksel
özellik çıkarılmış ve üç makine öğrenmesi modeli karşılaştırılmıştır. En iyi
genel sonuç Gradient Boosting modeliyle elde edilmiştir.

Sonuçlar, top1--top2 fark dizisinin doğru ve yanlış cevapları tek başına
kusursuz biçimde ayırmadığını; ancak bu diziden çıkarılan istatistiksel
özelliklerin cevap doğruluğu hakkında anlamlı ve ölçülebilir bir sinyal
taşıdığını göstermektedir. Bu nedenle top1--top2 fark analizi, özellikle dış
doğrulayıcı kullanılmadığı durumlarda LLM cevap güvenilirliği için yararlı bir
yardımcı sinyal olarak değerlendirilebilir.

\begin{thebibliography}{9}

\bibitem{cosmos}
YTU CE Cosmos, ``Turkish-Gemma-9b-T1,'' Hugging Face, 2024.
[Online]. Available: \url{https://huggingface.co/ytu-ce-cosmos/Turkish-Gemma-9b-T1}

\bibitem{gsm8k_tr}
YTU CE Cosmos, ``gsm8k\_tr,'' Hugging Face, 2024.
[Online]. Available: \url{https://huggingface.co/datasets/ytu-ce-cosmos/gsm8k_tr}

\bibitem{gsm8k}
K. Cobbe et al., ``Training Verifiers to Solve Math Word Problems,''
\emph{arXiv:2110.14168}, 2021.

\bibitem{hendrycks2016}
D. Hendrycks and K. Gimpel, ``A Baseline for Detecting Misclassified and
Out-of-Distribution Examples in Neural Networks,'' \emph{arXiv:1610.02136}, 2016.

\bibitem{kadavath2022}
S. Kadavath et al., ``Language Models (Mostly) Know What They Know,''
\emph{arXiv:2207.05221}, 2022.

\bibitem{selfcheckgpt}
P. Manakul, A. Liusie, and M. J. F. Gales, ``SelfCheckGPT: Zero-Resource
Black-Box Hallucination Detection for Generative Large Language Models,''
\emph{arXiv:2303.08896}, 2023.

\bibitem{sklearn}
F. Pedregosa et al., ``Scikit-learn: Machine Learning in Python,''
\emph{Journal of Machine Learning Research}, vol. 12, pp. 2825--2830, 2011.

\end{thebibliography}

\end{document}
